% GENERATED FILE. Edit canonical lessons, then run npm run generate:books.
% canonical-source-hash: fnv1a64:7cdf506eb1cce823
% canonical-lessons: JA-C70-mochi, JA-C70-nattou, JA-C70-chikuwa, JA-C70-hiyoko, JA-C70-hane

\chapter{Rice cake, Natto, Fish-paste tube, Chick, Wing}
\label{ch:ja-rice-cake-natto-fish-paste-tube-chick-wing}
\hlchaptermodality{70}

\begin{chapteropening}
\textbf{By the end of this chapter:} \emph{I can say a rice cake, natto, fermented soybeans, a fish-paste tube, a chick and a wing.}

Complete the last lesson of chapter 70: I can say a rice cake, natto, fermented soybeans, a fish-paste tube, a chick and a wing.
\end{chapteropening}

\section[\texorpdfstring{\ja{もち} (mochi)}{もち (mochi)}]{\ja{もち} (mochi) — a rice cake}

\label{lesson:JA-C70-mochi}



\begin{quote}
Before the new one: say the Japanese for ice, then the Japanese for the summer holidays.
\end{quote}

\subsection*{You'll want to know: \ja{もち}}
\textbf{\ja{もち}} — \emph{mochi} — \textquotedblleft{}a rice cake\textquotedblright{}.

\begin{cousinweb}
\textbf{\ja{もち}} (\emph{mochi}) is a native Japanese word, inherited rather than borrowed from Chinese, and its older origin is not settled. It is usually written in kana; its kanji, when used, is a kun'yomi attached later (JMdict, KANJIDIC2).
\end{cousinweb}

\begin{grammarlens}[title={What you've built}]
One more word for this chapter.
\end{grammarlens}

\subsection*{Your turn}
\begin{itemize}
\raggedright
  \item \emph{Say it:} \emph{mochi}
  \item \emph{Say it:} \emph{mochi}, once more
  \item \emph{Say it:} say \emph{mochi}
  \item \emph{Recall:} say the Japanese for ice, then the Japanese for the summer holidays, then say \emph{mochi} again
\end{itemize}

\begin{tcolorbox}[breakable,colback=teal!4,colframe=teal!35!black,title={Before you move on}]
What does \ja{もち} mean? (\textquotedblleft{}A rice cake\textquotedblright{}.) Say it once more.
\end{tcolorbox}
\section[\texorpdfstring{\ja{なっとう} (nattou)}{なっとう (nattou)}]{\ja{なっとう} (nattou) — natto, fermented soybeans}

\label{lesson:JA-C70-nattou}



\begin{quote}
Before the new one: say the Japanese for the summer holidays, then the Japanese for a rice cake.
\end{quote}

\subsection*{You'll want to know: \ja{なっとう}}
\textbf{\ja{なっとう}} — \emph{nattou} — \textquotedblleft{}natto, fermented soybeans\textquotedblright{}.

\begin{cousinweb}
\textbf{\ja{なっとう}} (\emph{nattou}) is Sino-Japanese (kango): two Chinese characters read with their on'yomi, readings borrowed from Chinese. They are \emph{nō}, \textquotedblleft{}receive, store\textquotedblright{} (Middle Chinese \emph{nop}) and \emph{tō}, \textquotedblleft{}bean\textquotedblright{} (\emph{duwH}) (Unihan, KANJIDIC2; Middle Chinese in Baxter's transcription). \emph{Nō} is shortened to \emph{nat} before \emph{tō}.
\end{cousinweb}

\begin{grammarlens}[title={What you've built}]
One more word for this chapter.
\end{grammarlens}

\subsection*{Your turn}
\begin{itemize}
\raggedright
  \item \emph{Say it:} \emph{nattou}
  \item \emph{Say it:} \emph{nattou}, once more
  \item \emph{Say it:} say \emph{nattou}
  \item \emph{Recall:} say the Japanese for the summer holidays, then the Japanese for a rice cake, then say \emph{nattou} again
\end{itemize}

\begin{tcolorbox}[breakable,colback=teal!4,colframe=teal!35!black,title={Before you move on}]
What does \ja{なっとう} mean? (\textquotedblleft{}Natto, fermented soybeans\textquotedblright{}.) Say it once more.
\end{tcolorbox}
\section[\texorpdfstring{\ja{ちくわ} (chikuwa)}{ちくわ (chikuwa)}]{\ja{ちくわ} (chikuwa) — a fish-paste tube}

\label{lesson:JA-C70-chikuwa}



\begin{quote}
Before the new one: say the Japanese for a rice cake, then the Japanese for natto.
\end{quote}

\subsection*{You'll want to know: \ja{ちくわ}}
\textbf{\ja{ちくわ}} — \emph{chikuwa} — \textquotedblleft{}a fish-paste tube\textquotedblright{}.

\begin{cousinweb}
\textbf{\ja{ちくわ}} (\emph{chikuwa}) mixes the two layers: \emph{chiku}, \textquotedblleft{}bamboo\textquotedblright{}, the on'yomi of its first character, borrowed from Chinese (Middle Chinese \emph{trjuwk}), and native \emph{wa}, \textquotedblleft{}ring\textquotedblright{}, read as a kun'yomi: a bamboo ring, from its shape (JMdict, KANJIDIC2).
\end{cousinweb}

\begin{grammarlens}[title={What you've built}]
One more word for this chapter.
\end{grammarlens}

\subsection*{Your turn}
\begin{itemize}
\raggedright
  \item \emph{Say it:} \emph{chikuwa}
  \item \emph{Say it:} \emph{chikuwa}, once more
  \item \emph{Say it:} say \emph{chikuwa}
  \item \emph{Recall:} say the Japanese for a rice cake, then the Japanese for natto, then say \emph{chikuwa} again
\end{itemize}

\begin{tcolorbox}[breakable,colback=teal!4,colframe=teal!35!black,title={Before you move on}]
What does \ja{ちくわ} mean? (\textquotedblleft{}A fish-paste tube\textquotedblright{}.) Say it once more.
\end{tcolorbox}
\section[\texorpdfstring{\ja{ひよこ} (hiyoko)}{ひよこ (hiyoko)}]{\ja{ひよこ} (hiyoko) — a chick}

\label{lesson:JA-C70-hiyoko}



\begin{quote}
Before the new one: say the Japanese for natto, then the Japanese for a fish-paste tube.
\end{quote}

\subsection*{You'll want to know: \ja{ひよこ}}
\textbf{\ja{ひよこ}} — \emph{hiyoko} — \textquotedblleft{}a chick\textquotedblright{}.

\begin{cousinweb}
\textbf{\ja{ひよこ}} (\emph{hiyoko}) is a native Japanese word, inherited rather than borrowed from Chinese, and its older origin is not settled. It is usually written in kana; its kanji, when used, is a kun'yomi attached later (JMdict, KANJIDIC2).
\end{cousinweb}

\begin{grammarlens}[title={What you've built}]
One more word for this chapter.
\end{grammarlens}

\subsection*{Your turn}
\begin{itemize}
\raggedright
  \item \emph{Say it:} \emph{hiyoko}
  \item \emph{Say it:} \emph{hiyoko}, once more
  \item \emph{Say it:} say \emph{hiyoko}
  \item \emph{Recall:} say the Japanese for natto, then the Japanese for a fish-paste tube, then say \emph{hiyoko} again
\end{itemize}

\begin{tcolorbox}[breakable,colback=teal!4,colframe=teal!35!black,title={Before you move on}]
What does \ja{ひよこ} mean? (\textquotedblleft{}A chick\textquotedblright{}.) Say it once more.
\end{tcolorbox}
\section[\texorpdfstring{\ja{はね} (hane)}{はね (hane)}]{\ja{はね} (hane) — a wing; a feather}

\label{lesson:JA-C70-hane}



\begin{quote}
Before the new one: say the Japanese for a fish-paste tube, then the Japanese for a chick.
\end{quote}

\subsection*{You'll want to know: \ja{はね}}
\textbf{\ja{はね}} — \emph{hane} — \textquotedblleft{}a wing; a feather\textquotedblright{}.

\begin{cousinweb}
\textbf{\ja{はね}} (\emph{hane}) is a native Japanese word, inherited rather than borrowed from Chinese, and its older origin is not settled. Its kanji was attached to it later as a kun'yomi, a native reading given to a Chinese character (JMdict, KANJIDIC2).
\end{cousinweb}

\begin{grammarlens}[title={What you've built}]
One more word for this chapter.
\end{grammarlens}

\subsection*{Your turn}
\begin{itemize}
\raggedright
  \item \emph{Say it:} \emph{hane}
  \item \emph{Say it:} \emph{hane}, once more
  \item \emph{Say it:} say \emph{hane}
  \item \emph{Recall:} say the Japanese for a fish-paste tube, then the Japanese for a chick, then say \emph{hane} again
\end{itemize}

\begin{tcolorbox}[breakable,colback=teal!4,colframe=teal!35!black,title={Before you move on}]
What does \ja{はね} mean? (\textquotedblleft{}A wing; a feather\textquotedblright{}.) Say it once more.
\end{tcolorbox}
